% =====================================================================
% 9. Comparison with the Basel IRB approach -- question (x)
% =====================================================================
\section{Comparison with the Basel IRB Approach (Question x)}
\label{sec:irb}

\subsection{REGULATORY CORRELATION, CAPITAL CHARGE AND MATURITY ADJUSTMENT}
For each counterparty we evaluate the formulas of Section~6.4 of the assignment with $\mathrm{AVC}=1$ and $M=1$:
\begin{align*}
\rho_k&=0.12\,w_k+0.24\,(1-w_k),\qquad w_k=\frac{1-e^{-50\pi_k}}{1-e^{-50}},\\
K_k&=(1-R_k)\left[\Phi\!\left(\frac{\Phi^{-1}(\pi_k)+\sqrt{\rho_k}\,\Phi^{-1}(0.999)}{\sqrt{1-\rho_k}}\right)-\pi_k\right]\mathrm{MA}(\pi_k,M).
\end{align*}
Table~\ref{tab:irb} reports the results by rating class. The regulatory asset correlation $\rho_k$ decreases from $0.24$ for the best ratings to $0.12$ for CCC, and the capital charge $K_k$ rises from \val{q10.Kbar_pct.AAA|f2}\% of exposure for AAA to \val{q10.Kbar_pct.CCC|f2}\% for CCC.

\begin{table}[tb]
\caption{Basel IRB quantities by rating class ($\mathrm{AVC}=1$, $M=1$). $\bar K_k$ and $\overline{\mathrm{LGD}}$ are unweighted averages within the class; amounts in USD. The last column reports the average systematic weight $\lambda_k$ of the model for comparison.}
\label{tab:irb}
\centering
\qrmtable{tab_irb}
\end{table}

\paragraph{Meaning of the maturity adjustment with $M=1$.}
Substituting $M=1$ into $\mathrm{MA}(\pi_k,M)=\{1+(M-2.5)\,b(\pi_k)\}/\{1-1.5\,b(\pi_k)\}$,
\[
\mathrm{MA}(\pi_k,1)=\frac{1+(1-2.5)\,b(\pi_k)}{1-1.5\,b(\pi_k)}=\frac{1-1.5\,b(\pi_k)}{1-1.5\,b(\pi_k)}=1
\]
for every $\pi_k$. The core of the IRB formula measures default risk over the one-year horizon of the default probability. A loan with a longer maturity can also lose value within the year without defaulting, through a rating downgrade, and the maturity adjustment adds capital for this migration risk; the slope $b(\pi_k)$ is larger for low default probabilities, because good credits have more room to deteriorate. With an effective maturity of exactly one year the maturity coincides with the horizon, there is no residual life after the horizon whose value could be impaired, and the adjustment is neutral: $K_k$ is a pure default-mode capital charge.

\paragraph{Portfolio capital.}
Summing over counterparties,
\[
K=\sum_{k=1}^{100}E_kK_k=\text{USD }\val{q10.K|usd},\quad \mathrm{RWA}=12.5\,K=\text{USD }\val{q10.RWA|usd} .
\]
$K$ is \val{q10.K_pct|pct2}\% of total exposure. The classes B and CCC account for \val{derived.K_share_B_CCC|pct0}\% of it.

\subsection{WHY THE IRB CAPITAL DIFFERS FROM THE MODEL-BASED UNEXPECTED LOSS}
\begin{table}[tb]
\caption{From the IRB capital to the model-based unexpected loss at $99.9\%$ (reference run).}
\label{tab:gap}
\centering
\qrmtable{tab_irb_gap}
\end{table}

\begin{figure}[tb]
\centering
\qrmfigure{fig_irb}
\caption{Left: regulatory asset correlation against the systematic weight $\lambda_k$ of the model, one point per counterparty (colour: default probability; line: equality). Right: unexpected loss at $99.9\%$ under the specifications of Table~\ref{tab:gap}.}
\label{fig:irb}
\end{figure}

The unexpected loss is USD~\val{q8.M2.UL_ref|usd} under M2 and USD~\val{q8.M3.UL_ref|usd} under M3, that is \val{q10.ratio_UL_M2|f2} and \val{q10.ratio_UL_M3|f2} times the IRB capital (under M1: \val{q10.ratio_UL_M1|f2} times). The IRB formula is the $99.9\%$ unexpected loss of a specific model: the asymptotic single-risk-factor model of \citet{Vasicek02} and \citet{Gordy03}, in which one Gaussian factor drives all obligors, each obligor has the prescribed correlation $\rho_k$, and the portfolio is infinitely granular, so that $L=\E[L\mid\text{factor}]$ and VaR is the sum of the conditional expected losses at the $99.9\%$ quantile of the factor. Our models differ from it in several respects, which Table~\ref{tab:gap} separates by changing one assumption at a time.
\begin{enumerate}
\item \emph{Level of the correlations.} In the IRB formula $\sqrt{\rho_k}$ is the loading of the obligor on the systematic factor; by (\ref{eq:assetcorr}) this role is played in our model by $\sqrt{\lambda_k}$. The data have an average $\lambda_k$ of \val{q10.lambda_mean|f2}, with values up to \val{q5.lam_max|f2}, against regulatory correlations between $0.12$ and $0.24$ (Figure~\ref{fig:irb}, left). Evaluating the IRB formula with $\rho_k=\lambda_k$ gives USD~\val{q10.K_lambda|usd}, \val{derived.K_lambda_ratio|f2} times $K$. This single change accounts for almost the entire gap.
\item \emph{Number of factors.} With two imperfectly correlated factors and heterogeneous loadings, obligors do not all move with the same systematic variable, which diversifies part of the systematic risk. The effect is small here, because by (\ref{eq:assetcorr}) what matters is the correlation between the combinations $a_j'\Theta$ and $a_k'\Theta$, which averages \val{q10.sys_corr_mean|f2} across pairs: with non-negative loadings on two indices that are themselves correlated at \val{q4.corr_all|f2}, the model is close to a one-factor model. The large-pool unexpected loss under M2 is USD~\val{q8.M2.UL_largepool|usd}, slightly below the one-factor value of the previous step.
\item \emph{Granularity.} With $100$ names, idiosyncratic risk is not fully diversified. Moving from the large pool to the actual portfolio adds USD~\val{q10.granularity_addon_M2|usd} under M2, \val{derived.gran_share_UL_M2|pct0}\% of the unexpected loss. Name concentration is thus a minor source of the gap: the largest possible single loss is USD~\val{q5.max_single_loss|usd}, small relative to a tail loss in which more than twenty names default together.
\item \emph{Distribution of the factors.} The IRB factor is Gaussian. A $t$-copula between the two factors (M3) adds \val{derived.UL_M3_vs_M2|pct1}\% to the unexpected loss; the empirical distribution (M1) adds \val{derived.UL_M1_vs_M2|pct0}\%.
\item \emph{Parameters prescribed, not estimated.} The regulatory $\rho_k$ is a function of the default probability alone, calibrated by the Basel Committee on portfolios of large internationally active banks \citep{BCBS05}. The $\lambda_k$ of the synthetic portfolio do not decrease with the default probability (last column of Table~\ref{tab:irb}), unlike the regulatory correlations.
\item \emph{Horizon.} The IRB formula refers to a one-year horizon. Our model applies the same $\pi_k$ to a one-week horizon with weekly factor returns. As shown in Section~\ref{sec:model}, after calibration of $d_k$ the horizon affects the results only through the shape of the factor distribution, so this difference is immaterial under M2, whose standardised factors are Gaussian at any horizon, and nearly so under M3, but not under M1, where the shape of weekly returns matters.
\end{enumerate}

\subsection{WHY A REGULATOR MAY PREFER THE SIMPLER FORMULA}
The IRB formula has properties that a portfolio model lacks and that matter for regulation.
\begin{itemize}
\item \emph{Portfolio invariance.} Under the single-factor, infinitely granular assumptions the capital for a loan depends only on its own characteristics and not on the rest of the portfolio \citep{Gordy03}. Capital is additive, $K=\sum_kE_kK_k$, can be computed loan by loan, and translates directly into risk weights. In a portfolio model the contribution of a loan depends on every other position.
\item \emph{Comparability and level playing field.} Two banks holding the same loan hold the same capital against it. Model-based capital would differ across banks for reasons unrelated to risk, such as the choice between M1, M2 and M3, which in our case moves the unexpected loss by about \val{derived.UL_M1_vs_M2|pct0}\%.
\item \emph{Limited scope for model risk and capital arbitrage.} Banks estimate the default probability and, in the advanced approach, the loss given default and the exposure, which can be validated against observed defaults and recoveries. Correlations and tail shapes, the parameters to which the tail is most sensitive and which are hardest to estimate and to backtest at the $99.9\%$ level, are fixed by the regulator and cannot be chosen so as to minimise capital.
\item \emph{Verifiability.} A closed-form formula can be checked by a supervisor line by line. A simulation-based figure depends on implementation choices and, as Section~\ref{sec:results} shows, on Monte Carlo error.
\end{itemize}
The price of this simplicity is that the formula ignores what is specific to a portfolio: name and sector concentration, and correlations that differ from the prescribed ones. In the Basel framework these are left to the supervisory review under Pillar~2. Our portfolio shows that the omission need not be conservative: for a portfolio with strong systematic exposure the IRB capital is about a third of the model-based unexpected loss (\val{q10.ratio_UL_M2|inv_pct0}\% under M2).
